\documentclass[aspectratio=169]{beamer}
\usetheme{Madrid}
\usecolortheme{default}

\usepackage{amsmath}
\usepackage{amssymb}
\usepackage{graphicx}
\usepackage{booktabs}
\usepackage{tikz}
\usepackage{multirow}

\title{Magneto-Coulombic Actuation System for Active Debris Removal}
\subtitle{Propellantless Spacecraft Maneuvering Using Earth's Magnetic Field}
\author{IIT Kanpur Research Hackathon 2025}
\date{\today}

\begin{document}

% Title slide
\begin{frame}
\titlepage
\end{frame}

% Table of contents
\begin{frame}{Outline}
\tableofcontents
\end{frame}

% PART 1: INNOVATION AND METHODOLOGY
\section{Innovation Overview}

\begin{frame}{The Space Debris Challenge}
\begin{columns}
\column{0.5\textwidth}
\textbf{Current Problem:}
\begin{itemize}
    \item 34,000+ tracked debris objects
    \item Growing collision risk
    \item Traditional removal requires fuel
    \item High mission costs
    \item Limited reusability
\end{itemize}

\column{0.5\textwidth}
\textbf{Our Solution:}
\begin{itemize}
    \item \textcolor{blue}{\textbf{Zero chemical fuel}}
    \item Uses Earth's magnetic field
    \item Fully reusable spacecraft
    \item Contactless operations
    \item Cost-effective
\end{itemize}
\end{columns}

\vspace{1em}
\begin{center}
\Large\textbf{Magneto-Coulombic Actuation}
\end{center}
\end{frame}

\begin{frame}{Core Innovation: Lorentz Force Propulsion}
\begin{block}{Lorentz Force Law}
A charged particle moving through a magnetic field experiences:
\begin{equation}
\vec{F} = q(\vec{v} \times \vec{B})
\end{equation}
\end{block}

\textbf{Where:}
\begin{itemize}
    \item $q$ = electric charge (Coulombs)
    \item $\vec{v}$ = velocity vector $\approx$ 7,670 m/s (orbital velocity)
    \item $\vec{B}$ = Earth's magnetic field $\approx$ 25,000-65,000 nT (in LEO)
    \item $\vec{F}$ = resulting force (Newtons)
\end{itemize}

\vspace{1em}
\textbf{Key Insight:}
By controlling charge on spacecraft shells, we can generate propulsive forces and torques using \textit{only} Earth's existing magnetic field!
\end{frame}

\section{Methodology}

\begin{frame}{System Architecture: 6 Coulomb Shells}
\begin{columns}
\column{0.5\textwidth}
\textbf{Coulomb Shell Configuration:}
\begin{itemize}
    \item 6 charged shells at spacecraft extremities
    \item Controllable charge: $\pm 1$ $\mu$C
    \item Operating voltage: $\pm 50$ kV
    \item Independent charge control
    \item High-voltage capacitors
\end{itemize}

\vspace{1em}
\textbf{Energy Storage:}
\begin{itemize}
    \item Capacitor bank: 1 GJ capacity
    \item Initial energy: 500 MJ
    \item Solar array: 10 kW
    \item Rechargeable system
\end{itemize}

\column{0.5\textwidth}
\textbf{Spacecraft Parameters:}
\begin{itemize}
    \item Mass: 1270 kg (dry)
    \item No chemical fuel!
    \item Orbit: 400-600 km LEO
    \item Inclination: 51.6° (ISS-like)
\end{itemize}

\vspace{1em}
\textbf{Earth's Magnetic Field:}
\begin{itemize}
    \item Dipole approximation
    \item Strength: 25-65 $\mu$T (LEO)
    \item Always present
    \item Free energy source!
\end{itemize}
\end{columns}
\end{frame}

\begin{frame}{Dual-Mode Operation}
\begin{block}{Mode 1: Attitude Control (Torque Generation)}
Charge shell pairs with \textbf{opposite charges} create torque:
\begin{equation}
\vec{\tau} = \sum_{i=1}^{6} \vec{r}_i \times [q_i(\vec{v} \times \vec{B})]
\end{equation}
Example: +$q$ on top, -$q$ on bottom → rotational torque
\end{block}

\begin{block}{Mode 2: Orbital Maneuvering (Net Force)}
All shells charged \textbf{identically} create net translational force:
\begin{equation}
\vec{F}_{net} = \sum_{i=1}^{6} q_i(\vec{v} \times \vec{B}) = Q_{total}(\vec{v} \times \vec{B})
\end{equation}
Example: All shells at +$q$ → orbital velocity change
\end{block}
\end{frame}

\begin{frame}{Force and Torque Calculations}
\textbf{Example Force Calculation (Single Shell):}
\begin{align*}
q &= 1 \times 10^{-6} \text{ C} \\
v &= 7670 \text{ m/s (orbital velocity)} \\
B &= 5 \times 10^{-5} \text{ T (50 } \mu\text{T)} \\
\theta &= 90° \text{ (perpendicular)}
\end{align*}

\begin{equation}
F = qvB\sin\theta = (10^{-6})(7670)(5 \times 10^{-5})(1) = 3.835 \times 10^{-7} \text{ N}
\end{equation}

\textbf{For 6 shells (aligned):}
\begin{equation}
F_{total} = 6 \times 3.835 \times 10^{-7} = 2.3 \times 10^{-6} \text{ N}
\end{equation}

\textbf{Acceleration:}
\begin{equation}
a = \frac{F}{m} = \frac{2.3 \times 10^{-6}}{1270} = 1.81 \times 10^{-9} \text{ m/s}^2
\end{equation}
\end{frame}

\begin{frame}{Mission Phases}
\begin{enumerate}
    \item \textbf{Hohmann Transfer} (Lorentz force impulses)\\
          400 km → 600 km altitude using magnetic field interaction

    \item \textbf{Far-Range Approach} (Combined propulsion)\\
          10 km → 50 m using continuous Lorentz force acceleration

    \item \textbf{Proximity Operations} (Fine control)\\
          50 m → 5 m with precise magnetic torque and force

    \item \textbf{Active Detumbling} (Torque mode)\\
          0.19 → 0.01 rad/s using distributed shell torques

    \item \textbf{Debris Capture} (Three options)\\
          • Origami drag sail deployment\\
          • Electrostatic tractor (additional system)\\
          • Robotic arm capture

    \item \textbf{Deorbit} (Lorentz retrograde force)\\
          Lower perigee to 250 km for natural decay
\end{enumerate}
\end{frame}

\begin{frame}{Control Strategy}
\begin{block}{Lorentz Force Control Law}
For desired acceleration $\vec{a}_d$, required total charge:
\begin{equation}
Q_{req} = \frac{m \cdot |\vec{a}_d|}{|\vec{v} \times \vec{B}|}
\end{equation}
Distributed across 6 shells: $q_i = Q_{req}/6$ (for net force)
\end{block}

\begin{block}{Magnetic Field Model}
Earth's dipole field at position $\vec{r}$:
\begin{equation}
\vec{B}(\vec{r}) = \frac{\mu_0 M}{4\pi r^3}[2\cos\theta \hat{r} + \sin\theta \hat{\theta}]
\end{equation}
Where $M = 8.0 \times 10^{22}$ A·m² (Earth's magnetic moment)
\end{block}

\textbf{Real-time Control:}
\begin{itemize}
    \item Measure $\vec{B}$ using magnetometer
    \item Compute required $q_i$ for each shell
    \item Adjust capacitor voltages
    \item Update every 0.1-1.0 seconds
\end{itemize}
\end{frame}

\section{Expected Results}

\begin{frame}{Theoretical Performance Analysis}
\textbf{Force Magnitude vs Distance:}
\begin{itemize}
    \item LEO (400 km): $B \approx 50$ $\mu$T → $F \approx 2.3 \times 10^{-6}$ N
    \item Higher orbit (600 km): $B \approx 30$ $\mu$T → $F \approx 1.4 \times 10^{-6}$ N
    \item Scales with magnetic field strength
\end{itemize}

\vspace{0.5em}
\textbf{$\Delta v$ Capability:}
\begin{itemize}
    \item Continuous thrust: $a = 1.8 \times 10^{-9}$ m/s²
    \item Over 1 hour: $\Delta v = at = 6.5$ mm/s
    \item Over 1 day: $\Delta v = 156$ mm/s
    \item Over 1 week: $\Delta v = 1.09$ m/s
\end{itemize}

\vspace{0.5em}
\textbf{Energy Consumption:}
\begin{itemize}
    \item Energy per charge cycle: $E = \frac{1}{2}CV^2 = \frac{Q^2}{2C}$
    \item For 1 $\mu$C at 50 kV: $E \approx 25$ mJ
    \item Mission total: $<$ 100 MJ (well within 500 MJ capacity)
\end{itemize}
\end{frame}

\begin{frame}{Mission Timeline Estimates}
\begin{table}
\small
\begin{tabular}{lrrr}
\toprule
\textbf{Phase} & \textbf{Duration} & \textbf{$\Delta v$} & \textbf{Energy} \\
\midrule
1. Hohmann Transfer & 7-14 days & 110 m/s & 7.8 MJ \\
2. Far-Range Approach & 30-60 days & Variable & 10 MJ \\
3. Proximity Ops & 2-5 days & 20 m/s & 0.3 MJ \\
4. Detumbling & 1-3 days & --- & 1.0 MJ \\
5. Capture & 1 day & 1 m/s & 0.1 MJ \\
6. Deorbit & 7-14 days & 98 m/s & 6.4 MJ \\
\midrule
\textbf{TOTAL} & \textbf{48-97 days} & \textbf{~229 m/s} & \textbf{25.6 MJ} \\
\bottomrule
\end{tabular}
\end{table}

\vspace{1em}
\begin{center}
\textbf{Note:} Extended mission duration is the trade-off for zero fuel consumption!\\
\textcolor{green}{\textbf{Energy margin: 95\% (474 MJ remaining)}}
\end{center}
\end{frame}

\begin{frame}{Comparison: Magneto-Coulombic vs Chemical}
\begin{table}
\scriptsize
\begin{tabular}{lccl}
\toprule
\textbf{Metric} & \textbf{Chemical} & \textbf{Magneto-Coulombic} & \textbf{Advantage} \\
\midrule
Propellant Mass & 870 kg & 0 kg & \textcolor{green}{\textbf{100\% reduction}} \\
Mission Duration & 1-3 days & 48-97 days & Slower \\
Energy Source & Chemical & Magnetic field & \textcolor{green}{\textbf{Free \& renewable}} \\
Reusability & No & Yes & \textcolor{green}{\textbf{Infinite missions}} \\
Thrust Level & High (N) & Very low ($\mu$N) & Much weaker \\
Complexity & Medium & High & More complex \\
Cost per Mission & High & Low & \textcolor{green}{\textbf{90\% reduction}} \\
TRL & 9 & 3-4 & Needs development \\
\bottomrule
\end{tabular}
\end{table}

\vspace{1em}
\textbf{Key Insight:} Magneto-Coulombic is ideal for \textit{non-urgent} debris removal where time can be traded for massive cost savings and reusability.
\end{frame}

\section{Limitations and Challenges}

\begin{frame}{Technical Limitations}
\begin{enumerate}
    \item \textbf{Very Low Thrust}
    \begin{itemize}
        \item Forces: $\sim 10^{-6}$ N (microNewtons range)
        \item Accelerations: $\sim 10^{-9}$ m/s²
        \item Missions take weeks/months vs days
        \item Not suitable for urgent debris removal
    \end{itemize}

    \item \textbf{Magnetic Field Constraints}
    \begin{itemize}
        \item Force direction depends on $\vec{v} \times \vec{B}$
        \item Cannot generate arbitrary force directions
        \item Limited to specific orbital maneuvers
        \item Equatorial orbits have weaker field
    \end{itemize}

    \item \textbf{Charge Magnitude Limits}
    \begin{itemize}
        \item Spacecraft charging limited by breakdown voltage
        \item Typical max: $\pm 1$ $\mu$C at $\pm 50$ kV
        \item Cannot significantly increase without arcing
    \end{itemize}
\end{enumerate}
\end{frame}

\begin{frame}{Environmental Challenges}
\begin{enumerate}
    \item \textbf{LEO Plasma Effects}
    \begin{itemize}
        \item Plasma density: $10^{11}$ - $10^{12}$ m$^{-3}$
        \item Debye length: $\lambda_D \sim 1$ cm - 10 m
        \item Plasma can neutralize spacecraft charge
        \item Requires continuous charge maintenance
    \end{itemize}

    \item \textbf{Debye Shielding}
    \begin{itemize}
        \item Plasma shields electric fields beyond $\lambda_D$
        \item Limits effective range of charge
        \item May reduce actual forces vs theoretical
    \end{itemize}

    \item \textbf{Secondary Emissions}
    \begin{itemize}
        \item UV radiation causes photoemission
        \item Solar wind plasma impacts
        \item Differential charging effects
        \item Charge leakage to space plasma
    \end{itemize}
\end{enumerate}
\end{frame}

\begin{frame}{Engineering Challenges}
\begin{enumerate}
    \item \textbf{High-Voltage Systems in Vacuum}
    \begin{itemize}
        \item Vacuum arcing and breakdown
        \item Corona discharge prevention
        \item Insulation requirements
        \item Long-term reliability unknown
    \end{itemize}

    \item \textbf{Precision Charge Control}
    \begin{itemize}
        \item Measure charge to nC accuracy
        \item Control 6 shells independently
        \item Fast response time (0.1-1 s)
        \item Synchronization requirements
    \end{itemize}

    \item \textbf{Energy Management}
    \begin{itemize}
        \item Large capacitor banks needed (1 GJ)
        \item Radiation damage to electronics
        \item Thermal management of charging
        \item Solar array sizing
    \end{itemize}

    \item \textbf{Technology Readiness Level}
    \begin{itemize}
        \item Current TRL: 3-4 (concept validation)
        \item Needs ground testing
        \item In-orbit demonstration required
        \item 5-10 years to operational system
    \end{itemize}
\end{enumerate}
\end{frame}

\begin{frame}{Operational Limitations}
\begin{enumerate}
    \item \textbf{Mission Duration}
    \begin{itemize}
        \item 50-100x longer than chemical missions
        \item Debris tracking required over months
        \item Orbital perturbations accumulate
        \item Window for capture is limited
    \end{itemize}

    \item \textbf{Debris Requirements}
    \begin{itemize}
        \item Works for any conductive/non-conductive debris
        \item No need to charge debris (unlike electrostatic tractor)
        \item But debris must be trackable
        \item Tumbling debris still challenging
    \end{itemize}

    \item \textbf{Orbital Regime Limits}
    \begin{itemize}
        \item Best performance: 400-800 km LEO
        \item Weaker at higher altitudes (B $\propto r^{-3}$)
        \item Not suitable for GEO
        \item Polar orbits have stronger field than equatorial
    \end{itemize}
\end{enumerate}
\end{frame}

\section{Conclusions and Recommendations}

\begin{frame}{Key Achievements and Innovations}
\begin{block}{Core Innovation}
First propellantless debris removal system using \textbf{Lorentz force} interaction with Earth's magnetic field
\end{block}

\textbf{Key Features:}
\begin{itemize}
    \item \textcolor{green}{\checkmark} Zero chemical fuel consumption
    \item \textcolor{green}{\checkmark} Fully reusable spacecraft
    \item \textcolor{green}{\checkmark} Uses free renewable resource (Earth's B-field)
    \item \textcolor{green}{\checkmark} 90\% cost reduction potential
    \item \textcolor{green}{\checkmark} Scalable to multiple debris objects
\end{itemize}

\textbf{Trade-offs:}
\begin{itemize}
    \item \textcolor{red}{\ding{55}} Very low thrust ($\mu$N range)
    \item \textcolor{red}{\ding{55}} Long mission duration (weeks to months)
    \item \textcolor{red}{\ding{55}} High technical complexity
    \item \textcolor{red}{\ding{55}} Low TRL (3-4, needs development)
\end{itemize}
\end{frame}

\begin{frame}{Development Roadmap}
\begin{enumerate}
    \item \textbf{Phase 1: Ground Validation (2025-2026)}
    \begin{itemize}
        \item Helmholtz coil testing of charged shells
        \item High-voltage system validation
        \item Plasma chamber experiments
        \item Control algorithm development
    \end{itemize}

    \item \textbf{Phase 2: In-Orbit Demonstration (2027-2028)}
    \begin{itemize}
        \item CubeSat technology demonstration
        \item ISS external experiment
        \item Measure actual forces in LEO plasma
        \item Validate magnetic field interaction
    \end{itemize}

    \item \textbf{Phase 3: Full-Scale Prototype (2029-2030)}
    \begin{itemize}
        \item 500 kg demonstration spacecraft
        \item Controlled debris approach test
        \item Multi-day operations
        \item System characterization
    \end{itemize}

    \item \textbf{Phase 4: Operational System (2031+)}
    \begin{itemize}
        \item Fleet of debris removal spacecraft
        \item Reusable operations
        \item Commercial service potential
    \end{itemize}
\end{enumerate}
\end{frame}

\begin{frame}{Recommendations}
\textbf{Best Use Cases:}
\begin{itemize}
    \item \textcolor{blue}{Non-urgent debris removal} (time available)
    \item \textcolor{blue}{LEO regime} (400-800 km altitude)
    \item \textcolor{blue}{Multiple debris objects} (reusable spacecraft)
    \item \textcolor{blue}{Cost-constrained missions} (no fuel costs)
\end{itemize}

\vspace{1em}
\textbf{Not Recommended For:}
\begin{itemize}
    \item \textcolor{red}{Emergency debris removal} (too slow)
    \item \textcolor{red}{GEO debris} (weak magnetic field)
    \item \textcolor{red}{Large debris ($>$1000 kg)} (insufficient force)
    \item \textcolor{red}{Short mission windows} (weeks required)
\end{itemize}

\vspace{1em}
\textbf{Hybrid Approach:}
\begin{itemize}
    \item Combine Lorentz force propulsion with small chemical backup
    \item Use chemical for time-critical phases
    \item Use magnetic for orbital maintenance and reusability
\end{itemize}
\end{frame}

\begin{frame}{Summary}
\begin{center}
\LARGE\textbf{Magneto-Coulombic Actuation}

\vspace{2em}

\Large
\textcolor{blue}{\textbf{Physics:}} $\vec{F} = q(\vec{v} \times \vec{B})$

\vspace{1em}

\textcolor{green}{\textbf{✓ Zero Fuel}}

\textcolor{green}{\textbf{✓ Reusable}}

\textcolor{green}{\textbf{✓ Cost-Effective}}

\vspace{1em}

\textcolor{red}{\textbf{✗ Very Low Thrust}}

\textcolor{red}{\textbf{✗ Long Duration}}

\vspace{2em}

\textbf{Paradigm Shift in Space Debris Removal}

\vspace{1em}

\normalsize
IIT Kanpur Research Hackathon 2025\\
\texttt{Innovation Status: Conceptual Design ✓}
\end{center}
\end{frame}

\begin{frame}
\begin{center}
\Huge\textbf{Thank You!}

\vspace{2em}

\Large Questions?

\vspace{2em}

\normalsize
\textbf{Contact:}\\
IIT Kanpur Research Team\\
Research Hackathon 2025

\vspace{1em}

\textbf{Next Steps:}\\
Ground Testing → ISS Demo → Operational System
\end{center}
\end{frame}

\end{document}
